\chapter{Evaluación de clasificadores}
\label{cap:evaluacion}

\section{Matriz de confusión}

Dado un umbral $t$ sobre el puntaje $s(\bm x)$, se predice crisis si $s(\bm x)\ge t$. La matriz de confusión cuenta verdaderos positivos (VP), falsos positivos (FP), verdaderos negativos (VN) y falsos negativos (FN). A partir de ella:
\begin{align}
\text{Sensibilidad} = \TPR(t) &= \frac{VP}{VP+FN}, &
\text{Especificidad} = 1-\FPR(t) &= \frac{VN}{VN+FP},\\
\text{VPP (precisión)} &= \frac{VP}{VP+FP}, &
F_1 &= \frac{2\,\text{VPP}\cdot\text{Sens.}}{\text{VPP}+\text{Sens.}}.
\end{align}
Con ventanas de $\Delta = 4$~s, las falsas alarmas por hora son $FP / \big((VN+FP)\,\Delta/3600\big)$.

\section{Curva ROC y AUC}

La curva ROC es el conjunto de puntos $(\FPR(t), \TPR(t))$ al recorrer todos los umbrales \parencite{fawcett2006}. El área bajo la curva tiene una interpretación probabilística:
\begin{proposicion}
$\mathrm{AUC} = P\big(s(\bm x^+) > s(\bm x^-)\big)$, donde $\bm x^+$ y $\bm x^-$ son una ventana de crisis y una normal elegidas al azar e independientes (contando los empates con peso $1/2$).
\end{proposicion}
\pendiente{demostración (relación con el estadístico U de Mann-Whitney).}

Con clases muy desbalanceadas, la curva ROC puede ser engañosamente optimista porque el FPR se calcula sobre una clase enorme. Por eso también se reporta la curva precisión-recall y su área (AP) \parencite{saito2015}.

\section{Elección del umbral como problema de optimización}

\subsection{Índice de Youden}
Sin restricciones, se elige $t$ que maximiza $J(t) = \TPR(t) - \FPR(t)$, el punto de la ROC más alejado de la diagonal.

\subsection{Criterio de Neyman-Pearson y multiplicador de Lagrange}

En un detector clínico la restricción natural es un máximo de falsas alarmas. Se plantea
\begin{equation}
\max_t\ \TPR(t)\quad\text{sujeto a}\quad\FPR(t)\le\alpha,
\label{eq:np}
\end{equation}
con $\alpha = 1\cdot\Delta/3600\approx 0{,}0011$, equivalente a una falsa alarma por hora. El lagrangiano es
\begin{equation}
\mathcal{L}(t,\lambda) = \TPR(t) - \lambda\big(\FPR(t)-\alpha\big),\qquad\lambda\ge 0.
\end{equation}
Si se parametriza la curva por $u = \FPR$ y $\TPR = g(u)$ es cóncava (lo que se garantiza reemplazando la ROC por su envolvente cóncava), las condiciones KKT son:
\[
g'(u^\star) = \lambda^\star,\qquad \lambda^\star(u^\star-\alpha) = 0,\qquad u^\star\le\alpha,\qquad\lambda^\star\ge 0.
\]
Como $g$ es creciente, la restricción está activa ($u^\star = \alpha$) y \textbf{el multiplicador óptimo es la pendiente de la ROC en el punto de operación}.

\begin{teorema}[Lema de Neyman-Pearson, \cite{neyman1933}]
Entre todos los tests con $\FPR\le\alpha$, el de mayor $\TPR$ rechaza la hipótesis nula (``normal'') cuando el cociente de verosimilitudes $\Lambda(s) = p(s\mid\text{crisis})/p(s\mid\text{normal})$ supera un umbral $\eta$ elegido para que $\FPR = \alpha$.
\end{teorema}

Además, si el puntaje es monótono en $\Lambda$, la pendiente de la ROC en el umbral $t$ es exactamente $\Lambda(t) = p(t\mid 1)/p(t\mid 0)$ \parencite{fawcett2006}. Por lo tanto $\lambda^\star = \eta$: el multiplicador de Lagrange del problema~\eqref{eq:np} es el umbral del lema de Neyman-Pearson, y se interpreta como cuántas veces más probable es el puntaje bajo crisis que bajo normalidad en la frontera de decisión. También es el ``precio sombra'' de la restricción: $d\,\TPR^\star/d\alpha = \lambda^\star$ indica cuánta sensibilidad se gana por cada unidad adicional de FPR permitida.

\pendiente{demostración del lema; figura \texttt{evaluacion\_final}.}

\section{Métricas por evento}

\pendiente{sensibilidad por crisis y latencia de detección.}

\section{Pruebas estadísticas para comparar modelos}

\pendiente{Friedman, Wilcoxon de rangos con signo, corrección de Holm, ANOVA de dos factores con interacción \parencite{demsar2006}.}
